\clearpage
\VerdierPhysicalPageAnchor{0230}%
\begingroup\small\setlength{\parskip}{0pt}\setlength{\emergencystretch}{1.5em}
\setlength{\abovedisplayskip}{4pt plus 2pt minus 2pt}%
\setlength{\belowdisplayskip}{4pt plus 2pt minus 2pt}%

\VerdierSemanticAnchor{ch4d-4-1-6-cont-2}
\noindent 然后置

\VerdierSemanticAnchor{ch4d-4-1-6-6}
\begin{equation}
\int_n f\bigl((p,q),(p',q')\bigr)
=h\bigl((p,q),(p',q')\bigr).
\tag{4.1.6.6}
\end{equation}

并且令

\VerdierSemanticAnchor{ch4d-4-1-6-7}
\begin{equation}
{}^{m}\!d(p,q,r)\colon{}^{m}\!Y(q,r)
\longrightarrow{}^{m}\!Y(p,q)[1]
\tag{4.1.6.7}
\end{equation}

为复合态射

\VerdierSemanticAnchor{ch4d-4-1-6-8}
\begin{equation}
\begin{aligned}
{}^{m}\!Y(q,r)
&\xrightarrow{\int_n g(p,q,r)}
\int_n T_m\bigl({}^{m}\!X(p,q)\bigr)
\xrightarrow{(*)}
\left(\int_n{}^{m}\!X(p,q)\right)[1]\\
&=Y(p,q)[1].
\end{aligned}
\tag{4.1.6.8}
\end{equation}

其中态射 $(*)$ 是关联单复形函子与平移交换的同构（第一章，2.3.2.1）。

由此得到单复形 $Y=\int_n X$ 的一个滤过：

\VerdierSemanticAnchor{ch4d-4-1-6-9}
\begin{equation}
0={}^{m}\!Y(-\infty)\subset\cdots\subset{}^{m}\!Y(q)
\subset\cdots\subset{}^{m}\!Y(+\infty)=Y
\tag{4.1.6.9}
\end{equation}

称为\emph{由 $X$ 的第 $m$ 个次数定义的滤过}。

通过 (4.1.1) 中回顾的方法，可以把一个取值于 $K(\mathcal A)$ 的谱对象
关联给滤过复形 (4.1.6.9)，记作

\VerdierSemanticAnchor{ch4d-4-1-6-10}
\begin{equation}
\bigl(\mathcal Y(p,q),\delta(p,q,r)\bigr).
\tag{4.1.6.10}
\end{equation}

最后，在下面的命题中，我们将滥用记号，以同一记号表示
$\operatorname{comp}(\mathcal A)$ 的一个态射及其在 $K(\mathcal A)$ 中的像。

\VerdierSemanticAnchor{ch4d-4-1-7}
\medskip\noindent\textbf{命题 4.1.7.}\quad

\VerdierSemanticAnchor{ch4d-4-1-7-a}
\noindent\textit{a) 对象 ${}^{m}\!Y(p,q)$ (4.1.6.5)、态射
$h\bigl((p,q),(p',q')\bigr)$ (4.1.6.6) 以及态射
${}^{m}\!d(p,q,r)$ (4.1.6.7) 定义了一个取值于 $K(\mathcal A)$ 的谱对象，
它函子性地依赖于 $n$-重复形 $X$。}

\noindent\textit{这个谱对象称为由按第 $m$ 个次数滤过的复形 $X$ 所定义的、
取值于 $K(\mathcal A)$ 的谱对象。}

\VerdierSemanticAnchor{ch4d-4-1-7-b}
\noindent\textit{b) 谱对象
$\bigl(\mathcal Y(p,q),\delta(p,q,r)\bigr)$ (4.1.6.10) 与
$\bigl({}^{m}\!Y(p,q),{}^{m}\!d(p,q,r)\bigr)$ 之间存在典范同构（即关于
$\operatorname{comp}^n(\mathcal A)$ 的对象 $X$ 函子性的同构）。}
\endgroup
